\documentclass[11pt]{article}
\usepackage[margin=1in]{geometry}
\usepackage{amsmath,amssymb,amsthm}
\usepackage{mathtools}
\usepackage{enumitem}
\usepackage{hyperref}
\usepackage{cleveref}
\usepackage{booktabs}
\usepackage{tabularx}
\usepackage{microtype}
\usepackage{cite}
\usepackage{xurl}
\usepackage{path}
\hypersetup{colorlinks=true,linkcolor=blue,citecolor=blue,urlcolor=blue}
\newtheorem{definition}{Definition}
\newtheorem{lemma}{Lemma}
\newtheorem{remark}{Remark}
\title{Public Request Carryforward for Source-Only Successor Releases\\\large When Paper-Level Request Status Persists, Advances, or Reopens across Release Evolution}
\author{Anonymous}
\date{Unified archive note v0.08 (2026-03-21; archive v467)}
\begin{document}
\maketitle

\begin{abstract}
The archive can now say what a source-only paper promises on request, what omitted-support families that promise covers, when each family has been fulfilled, and what the overall paper-level public status token is for one release. One cross-release question still recurs: when a successor release ships, does that public request status remain valid, advance, or reopen? This note gives that carryforward rule one home. It defines compact \emph{public request carryforward envelopes}: successor-release maintained objects stating whether a predecessor paper-level source-only request status survives unchanged, advances, or reopens under the new release pair.
\end{abstract}

\paragraph{Non-redundancy note.}
Synthesis~32 owns the stable release spine for successor maintenance, Synthesis~56 owns exact paper-facing public request contracts, Synthesis~58 owns family-scoped completion criteria, Synthesis~59 owns the within-release paper-level public status token, Synthesis~61 owns the visible outward-facing carryforward wrapper sentence, Synthesis~73 owns the request-side terminal stop later terse papers should usually cite, Synthesis~74--75 own the paired terminal / cutover discipline, and Synthesis~17 owns the concrete worked instantiation.\cite{series:releasespine,series:publicrequestcontracts,series:requestfulfillmentcertificates,series:publicrequeststatusenvelopes,series:publicrequestnotices,series:publicrequestresponsepacketrefreshresponsepacketclosureverdicts,series:pairedterminalcitationdiscipline,series:pairedterminalcutoverrule,series:workedexample}
This note owns only the cross-release carryforward rule saying whether one already-public paper-level source-only status token persists, advances, or reopens under a successor release. It does not own the outward-facing sentence readers should see once that carryforward verdict is known, and later terse papers should name it directly only when that exact paper-level successor verdict is the live issue.

\paragraph{Scope.}
This is not a new release stage, not a new paper-facing contract, and not a replacement for the within-release public status envelope.
It is a thin successor-release bridge saying how one paper-level source-only status token behaves when a successor cut ships.

\paragraph{Exports.}
This note exports (i) public request carryforward envelopes, (ii) a status-preservation rule, (iii) a family-stability rule, (iv) an advance/reopen rule, (v) a shared-carryforward-token-vocabulary rule, (vi) a closure-import rule, and (vii) a successor-release compression rule saying that later papers may cite one compact maintained object when they need the cross-release state of a source-only request promise rather than restating it in prose. If they instead need the exact outward-facing sentence readers should see for that carryforward verdict, they should import Synthesis~61 rather than widen this note. If they need only the stable request-side stop or the stop-versus-widen judgment for the adjacent visible source-only branch, they should cite Synthesis~73--75 rather than restitch the wrapper ladder here.

\paragraph{Later-terse-paper citation posture.}
Later terse papers should cite Synthesis~60 only when the live issue is whether one already-public paper-level source-only status token persisted, advanced, or reopened across one concrete successor release. Stop at Synthesis~59 when the live issue is only which within-release paper-level public token is true. Stop at Synthesis~61--63 when the carryforward verdict is already fixed and the remaining issue is the exact outward-facing sentence, canonical notice family, or token-keyed selector that readers were shown. Widen through Synthesis~64--72 only when that visible wrapper layer is already fixed and the live issue has moved into the exact follow-up handoff, predecessor packet cut, successor reuse verdict, packet-local refresh delta, visible refresh sentence, canonical refresh family, family-selection judgment, successor-facing handoff, or exact successor-facing packet cut. Widen to Synthesis~73 when the live issue is whether that visible request-side branch already closes. Widen to Synthesis~74--75 when the remaining issue is paired-terminal reuse or fail-closed cutover. Stop at Synthesis~17 when the concrete maintained worked route is the live issue instead.\cite{series:publicrequeststatusenvelopes,series:publicrequestnotices,series:publicrequestnoticenormalforms,series:publicrequestnoticeselections,series:publicrequestresponsemenus,series:publicrequestresponsepackets,series:publicrequestresponsepacketcarryforward,series:publicrequestresponsepacketdeltaledgers,series:publicrequestresponsepacketrefreshnotices,series:publicrequestresponsepacketrefreshnoticenormalforms,series:publicrequestresponsepacketrefreshnoticeselections,series:publicrequestresponsepacketrefreshresponsemenus,series:publicrequestresponsepacketrefreshresponsepackets,series:publicrequestresponsepacketrefreshresponsepacketclosureverdicts,series:pairedterminalcitationdiscipline,series:pairedterminalcutoverrule,series:workedexample}

\section{Why carryforward deserves its own note}
Within one release, a public request status envelope answers the paper-level question ``is the source-only promise complete, partial, or open?''
Across releases, one more question appears.
If the successor release preserves the same promised families and their completion remains intact, the public status may carry forward unchanged.
If the successor release adds or changes promised families, the old complete token may have to reopen.
If the successor release closes an earlier gap, the status may advance.
Without one maintained cross-release object, those cases drift into prose and later notes silently disagree about whether the old promise survived the successor cut.

\begin{definition}[Public request carryforward envelope]
Fix a predecessor public request status envelope $E_0$, a successor public request contract $r_1$, a successor public request status envelope $E_1$, and the successor-release continuity/closure context from the maintained release spine.
A \emph{public request carryforward envelope} is a tuple
\[
C=(E_0,r_1,E_1,\chi,\pi,\Delta,\Gamma,\rho,\omega)
\]
containing imported continuity/closure labels, changed/reopened-family sets, one carryforward token $\rho\in\{\texttt{preserved\_complete},\texttt{preserved\_noncomplete},\texttt{advanced},\texttt{reopened}\}$, and one explanation field $\omega$.
The envelope does not replace release classification or within-release completion.
It only states how the public request status behaves across one concrete release transition.
\end{definition}

\begin{definition}[Status-preservation rule]
A carryforward envelope satisfies the \emph{status-preservation rule} if it emits \texttt{preserved\_complete} exactly when both predecessor and successor public status tokens are \texttt{complete} and $\Gamma=\varnothing$, and emits \texttt{preserved\_noncomplete} exactly when the predecessor and successor public status tokens agree on a non-complete status and $\Gamma=\varnothing$.
So preservation is a release-pair verdict about the public token together with the absence of reopened families, not a new within-release completion test.
\end{definition}

\begin{definition}[Family-stability rule]
A carryforward envelope satisfies the \emph{family-stability rule} if every key in $\Delta\cup\Gamma$ is already covered by the successor public request contract and if a family may appear in $\Gamma$ only when the successor status basis for that family is no longer preserved.
So the carryforward envelope may report where continuity broke, but it may not invent new family keys or silently reopen an unchanged one.
\end{definition}

\begin{definition}[Advance and reopen rules]
A carryforward envelope satisfies the \emph{advance rule} if it emits \texttt{advanced} exactly when the successor paper-level status is strictly stronger than the predecessor status under the order \texttt{open} $<$ \texttt{partial} $<$ \texttt{complete}.
It satisfies the \emph{reopen rule} if it emits \texttt{reopened} whenever a predecessor-complete promise becomes non-complete or whenever $\Gamma\neq\varnothing$.
So the carryforward owner is the cross-release verdict, not the within-release proof.
\end{definition}

\begin{definition}[Shared-carryforward-token-vocabulary rule]
A carryforward envelope satisfies the \emph{shared-carryforward-token-vocabulary rule} if its token $\rho$ may recur in many later notice, normal-form, selector, route, or spine objects only as imported cross-release vocabulary: the same carryforward token name may be reused across many visible wrappers, but the truth of that token remains release-pair-scoped to the concrete carryforward envelope that emits it.
So the token may travel, but the release-pair verdict remains anchored to the envelope.
\end{definition}

\begin{definition}[Closure-import rule]
A carryforward envelope satisfies the \emph{closure-import rule} if its continuity decision $\chi$ and publication-closure status $\pi$ are imported from already-maintained successor-release continuity/closure objects and are not redefined by the envelope itself.
So the envelope may summarize why a token was preserved, advanced, or reopened, but it may not become the new owner of release-pair continuity or publication closure.
\end{definition}

\begin{lemma}[Carryforward envelopes stabilize cross-release public request status without replacing release or completion owners]
If a successor release keeps the release-spine, public request contract, and within-release public status envelope owners intact, and if a companion carryforward envelope satisfies the status-preservation, family-stability, advance/reopen, shared-carryforward-token-vocabulary, and closure-import rules, then a later terse paper may cite one carryforward envelope to explain whether a source-only request promise survived, advanced, or reopened across the successor release without re-owning the narrower release or completion semantics.
\end{lemma}

\begin{proof}
The within-release public status envelope still owns the successor paper-level token, and the release spine plus publication-closure objects still own the successor maintenance context.
The carryforward envelope reads those maintained owners and adds only one new cross-release statement: whether the old public request status survives unchanged, improves, or reopens.
The family-stability rule prevents the envelope from inventing new family keys or silently reopening unchanged families, while the closure-import rule prevents it from taking ownership of the continuity and closure decisions it cites.
The shared-carryforward-token-vocabulary rule then makes explicit that the same carryforward token name may be imported into many later visible wrappers without moving cross-release ownership away from the concrete predecessor/successor envelope pair.
Because the advance/reopen rules are defined from imported predecessor/successor tokens and reopened-family information, the carryforward envelope cannot silently redefine within-release completion or release classification.
So it gives later papers one stable cross-release answer without widening any narrower owner.\qedhere
\end{proof}

\begin{table}[t]
\centering
\small
\begin{tabularx}{0.97\linewidth}{@{}p{0.35\linewidth}p{0.19\linewidth}X@{}}
\toprule
\textbf{Imported release-pair pattern} & \textbf{Forced carryforward token} & \textbf{Why / next owner}\\
\midrule
Predecessor and successor tokens are both \texttt{complete} and \texttt{reopened\_family\_keys = []} & \texttt{preserved\_complete} & Preservation is honest only when the successor keeps full paper-level completeness without reopening any covered family.\\
Predecessor and successor tokens agree on a non-complete value and \texttt{reopened\_family\_keys = []} & \texttt{preserved\_noncomplete} & The release pair preserves one already-public non-complete status without claiming new progress.\\
Successor token is strictly stronger than predecessor and \texttt{reopened\_family\_keys = []} & \texttt{advanced} & The release pair improved the paper-level source-only status rather than merely preserving it.\\
The reopened-family set is nonempty, or a predecessor-\texttt{complete} token weakens in the successor & \texttt{reopened} & Reopened families or lost completeness dominate any matching token labels; widen to Synthesis~61 only after this release-pair verdict is fixed.\\
\bottomrule
\end{tabularx}
\caption{The carryforward envelope owns only a release-pair compression of imported within-release tokens plus reopened-family evidence. It sits strictly between the within-release token layer and the later visible notice sentence.}
\label{tab:public-request-carryforward-token-matrix}
\end{table}

\paragraph{Stop discipline.}
Later terse papers should cite Synthesis~60 only when the live issue is whether one already-public paper-level source-only status token persisted, advanced, or reopened across one concrete successor release.
Stop at Synthesis~59 when the live issue is only which within-release paper-level public token is true.
Stop at Synthesis~61--63 when the carryforward verdict is already fixed and the remaining issue is the exact outward-facing sentence, canonical notice family, or token-keyed selector that readers were shown.
Widen through Synthesis~64--72 only when that visible wrapper layer is already fixed and the live issue has moved into the exact follow-up handoff, predecessor packet cut, successor reuse verdict, packet-local refresh delta, visible refresh sentence, canonical refresh family, family-selection judgment, successor-facing handoff, or exact successor-facing packet cut.
Widen to Synthesis~73 when the live issue is whether that visible request-side branch already closes.
Widen to Synthesis~74--75 when the remaining issue is paired-terminal reuse or fail-closed cutover.
Stop at Synthesis~17 when the concrete maintained worked route is the live issue instead.\cite{series:publicrequestresponsemenus,series:publicrequestresponsepackets,series:publicrequestresponsepacketsuccessorcarryforward,series:publicrequestresponsepacketdeltaledgers,series:publicrequestresponsepacketrefreshnotices,series:publicrequestresponsepacketrefreshnormalforms,series:publicrequestresponsepacketrefreshselections,series:publicrequestresponsepacketrefreshresponsemenus,series:publicrequestresponsepacketrefreshresponsepackets,series:publicrequestresponsepacketrefreshresponsepacketclosureverdicts,series:pairedterminalcitationdiscipline,series:pairedterminalcutoverrule}

\paragraph{A minimal public carryforward card.}
\begin{table}[t]
\centering
\small
\begin{tabularx}{0.97\linewidth}{@{}p{0.31\linewidth}X@{}}
\toprule
\textbf{Public row} & \textbf{Pinned value}\\
\midrule
Envelope anchor & \texttt{public\_request\_carryforward\_envelope\_key = worked\_example\_receipt\_interlock\_public\_request\_carryforward\_envelope} from \nolinkurl{example_public_request_carryforward_envelope.json} \\
Release pair & \texttt{base\_release\_id = worked-example-draft-194} and \texttt{successor\_release\_id = worked-example-draft-194b} \\
Imported basis floor & \texttt{predecessor\_public\_request\_status\_envelope\_key = worked\_example\_receipt\_interlock\_public\_request\_status\_envelope} and \texttt{successor\_public\_request\_contract\_key = worked\_example\_receipt\_interlock\_public\_request\_contract} \\
Imported token pair & \texttt{predecessor\_public\_request\_status\_token = complete} and \texttt{successor\_public\_request\_status\_token = complete} \\
Continuity / closure floor & \texttt{continuity\_decision = same\_certified\_interface\_replayed\_surfaces} and \texttt{publication\_closure\_status = complete} \\
Family drift floor & \texttt{changed\_family\_keys = []} and \texttt{reopened\_family\_keys = []} \\
Derived carryforward verdict & \texttt{carryforward\_token = preserved\_complete} with \texttt{carryforward\_token\_mode = shared\_vocabulary\_release\_pair\_scoped\_verdict} \\
Escalation pointer & Widen to Synthesis~59 for the within-release paper-level token, Synthesis~61--63 for the visible wrapper / canonical family / selector layer, and Synthesis~73--75 once the live issue is request-side terminal reuse or cutover \\
\bottomrule
\end{tabularx}
\caption{A minimal public carryforward card. This is the smallest public answer later terse papers can quote directly when the live issue is whether one already-public paper-level source-only status token persisted, advanced, or reopened across one concrete successor release rather than the narrower within-release token basis, the visible wrapper sentence, or the request-side terminal tail.}
\label{tab:minimal-public-request-carryforward-card}
\end{table}

This card is intentionally non-decorative: it pins one predecessor/successor release pair, one imported status-envelope/contract basis, one continuity/closure floor, one changed/reopened-family floor, and one release-pair-scoped carryforward verdict without silently re-owning the within-release token, the visible notice sentence, or the request-side terminal stop.\cite{series:workedexample}

\paragraph{One tiny carryforward drill.}
In the worked carryforward envelope, the predecessor and successor paper-level public request tokens are both \texttt{complete}, the continuity decision is \texttt{same\_certified\_interface\_replayed\_surfaces}, the publication-closure status stays \texttt{complete}, and both the changed-family and reopened-family sets are empty. By the preservation and reopen rules, that leaves exactly one honest cross-release verdict: \texttt{carryforward\_token = preserved\_complete}. If either covered family reopened or the successor paper-level token weakened, the same carryforward owner would have to publish \texttt{reopened} instead of \texttt{preserved\_complete}; if the successor paper-level token strengthened from \texttt{open} or \texttt{partial} to \texttt{complete}, the same owner would have to publish \texttt{advanced}. Later terse papers may therefore quote the carryforward token directly only as long as that imported predecessor/successor token pair and reopened-family floor remain unchanged.\cite{series:workedexample}

\section{Worked example pointer}
The maintained worked bundle now ships \nolinkurl{example_public_request_carryforward_envelope.json}.\cite{series:workedexample}
It imports the worked note's paper-level public request status envelope together with the canned successor-release comparison context and states whether the source-only public request status survives that successor cut. In the shipped worked envelope the continuity decision is \texttt{same\_certified\_interface\_replayed\_surfaces}, the publication-closure status is \texttt{complete}, the changed-family and reopened-family sets are both empty, and the predecessor/successor paper-level public tokens are both \texttt{complete}, so the carryforward token is \texttt{preserved\_complete}.
That gives later terse papers one citeable cross-release answer when they need to say that a source-only public request promise survived a successor release without reopening. The worked example now also exports \texttt{carryforward\_token\_mode=shared\_vocabulary\_release\_pair\_scoped\_verdict} on that envelope and mirrors the same carryforward-token vocabulary posture through its answer routes and series spine, so later bridge papers can stop at one shipped route object while still seeing that the token remains release-pair vocabulary rather than visible-wrapper ownership. Those later terse papers should usually combine that route-visible carryforward basis with Synthesis~73--75 when they need only the stable request-side stop or stop-versus-widen judgment, and widen through Synthesis~61--63 only when the exact outward-facing sentence, canonical family, or selector is itself disputed.\cite{series:publicrequestnotices,series:publicrequestnoticenormalforms,series:publicrequestnoticeselections,series:publicrequestresponsepacketrefreshresponsepacketclosureverdicts,series:pairedterminalcitationdiscipline,series:pairedterminalcutoverrule}

\begin{thebibliography}{99}\small

\bibitem{series:releasespine}
Anonymous.
\newblock Stable Release Spine Contracts for Successor Maintenance: Stage Ownership, Handoff Rules, and Auditable Walk-Throughs.
\newblock Series synthesis note (Synthesis~32), 2026. (This archive.)

\bibitem{series:publicrequestcontracts}
Anonymous.
\newblock Public Request Contracts for Source-Only Releases: Exact Paper-Level Handoffs from Boilerplate Policy to Maintained Packet Cuts.
\newblock Series synthesis note (Synthesis~56), 2026. (This archive.)

\bibitem{series:requestfulfillmentcertificates}
Anonymous.
\newblock Request Fulfillment Certificates for Source-Only Releases: Exact Completion Criteria for Public Request Contracts and Typed Omitted-Support Cuts.
\newblock Series synthesis note (Synthesis~58), 2026. (This archive.)

\bibitem{series:publicrequeststatusenvelopes}
Anonymous.
\newblock Public Request Status Envelopes for Source-Only Releases: Paper-Level Tokens Compressing Family Completion without Re-owning the Promise.
\newblock Series synthesis note (Synthesis~59), 2026. (This archive.)

\bibitem{series:workedexample}
Anonymous.
\newblock Worked Example: Receipt Interlock for an Anonymous-DHT Reader-Privacy Claim.
\newblock Series synthesis note (Synthesis~17), 2026. (This archive.)

\bibitem{series:publicrequestnotices}
Anonymous.
\newblock Public Request Notices for Source-Only Releases: Outward-Facing Sentences for Within-Release Status and Cross-Release Carryforward.
\newblock Series synthesis note (Synthesis~61), 2026. (This archive.)

\bibitem{series:publicrequestnoticenormalforms}
Anonymous.
\newblock Public Request Notice Normal Forms for Source-Only Releases: Canonical Sentence Families for Within-Release Status and Cross-Release Carryforward Notices.
\newblock Series synthesis note (Synthesis~62), 2026. (This archive.)

\bibitem{series:publicrequestnoticeselections}
Anonymous.
\newblock Public Request Notice Selections for Source-Only Releases: Choosing Canonical Notice Families from Within-Release Status and Cross-Release Carryforward Tokens.
\newblock Series synthesis note (Synthesis~63), 2026. (This archive.)

\bibitem{series:publicrequestresponsemenus}
Anonymous.
\newblock Public Request Response Menus for Source-Only Releases: Smallest-Sufficient Follow-Ups from Notice Families to Maintained Packet Cuts.
\newblock Series synthesis note (Synthesis~64), 2026. (This archive.)

\bibitem{series:publicrequestresponsepackets}
Anonymous.
\newblock Public Request Response Packets for Source-Only Releases: Exact Packet Cuts for Maintained Follow-Ups.
\newblock Series synthesis note (Synthesis~65), 2026. (This archive.)

\bibitem{series:publicrequestresponsepacketsuccessorcarryforward}
Anonymous.
\newblock Public Request Response Packet Carryforward for Source-Only Successor Releases: When Maintained Packet Cuts Persist, Refresh, or Reopen.
\newblock Series synthesis note (Synthesis~66), 2026. (This archive.)

\bibitem{series:publicrequestresponsepacketdeltaledgers}
Anonymous.
\newblock Public Request Response Packet Delta Ledgers for Source-Only Successor Releases: Exact Packet-Local Refresh Deltas under Carryforward.
\newblock Series synthesis note (Synthesis~67), 2026. (This archive.)

\bibitem{series:publicrequestresponsepacketrefreshnotices}
Anonymous.
\newblock Public Request Response Packet Refresh Notices for Source-Only Successor Releases: Visible Wrapper Sentences after Packet-Local Refresh.
\newblock Series synthesis note (Synthesis~68), 2026. (This archive.)

\bibitem{series:publicrequestresponsepacketrefreshnormalforms}
Anonymous.
\newblock Public Request Response Packet Refresh Notice Normal Forms for Source-Only Successor Releases: Canonical Visible Wrapper Families after Packet Refresh.
\newblock Series synthesis note (Synthesis~69), 2026. (This archive.)

\bibitem{series:publicrequestresponsepacketrefreshselections}
Anonymous.
\newblock Public Request Response Packet Refresh Notice Selections for Source-Only Successor Releases: Choosing Canonical Wrapper Families after Packet Refresh.
\newblock Series synthesis note (Synthesis~70), 2026. (This archive.)

\bibitem{series:publicrequestresponsepacketrefreshresponsemenus}
Anonymous.
\newblock Public Request Response Packet Refresh Response Menus for Source-Only Successor Releases: Smallest-Sufficient Successor Handoffs after Visible Refresh.
\newblock Series synthesis note (Synthesis~71), 2026. (This archive.)

\bibitem{series:publicrequestresponsepacketrefreshresponsepackets}
Anonymous.
\newblock Public Request Response Packet Refresh Response Packets for Source-Only Successor Releases: Exact Successor-Facing Packet Cuts after Visible Refresh.
\newblock Series synthesis note (Synthesis~72), 2026. (This archive.)

\bibitem{series:publicrequestresponsepacketrefreshresponsepacketclosureverdicts}
Anonymous.
\newblock Public Request Response Packet Refresh Response Packet Closure Verdicts for Source-Only Successor Releases: Capstone Certificates for Exact Successor-Facing Handoff Slices.
\newblock Series synthesis note (Synthesis~73), 2026. (This archive.)

\bibitem{series:pairedterminalcitationdiscipline}
Anonymous.
\newblock Paired Terminal Citation Discipline for Stable Review Spines: One Answer-Side Stop and One Request-Side Capstone for Terse Papers.
\newblock Series synthesis note (Synthesis~74), 2026. (This archive.)

\bibitem{series:pairedterminalcutoverrule}
Anonymous.
\newblock Paired Terminal Cutover Rules for Stable Review Spines: When Later Terse Papers Must Stop, Widen, or Reopen.
\newblock Series synthesis note (Synthesis~75), 2026. (This archive.)

\end{thebibliography}
\end{document}
